% !TeX root = ../papers/nc-p3.tex
\section{階数15の分岐上の全点正則性と種数五の曲線族}
\label{sec:rank-fifteen-all-ramification}
\begin{theorem}\label{thm:rank-fifteen-all-ramification}
標数零で二つの階数15分岐曲線の全ての通常窓は標準Koszul左右AS正則である。
各始点で道空間の直和は四つの中心四次式の多項式環上で階数256の自由加群となる。
通常倍率の開集合を、$zk^2=1$ と $\delta=0$ の両軌跡を全て含み、
その上の両四重被覆の全窓が同じ正則性を持つように取れる。
\end{theorem}
\begin{proof}
前節により、101で通常な次数2以上の剰余における長さ13の商だけが残る。
$S=\mathbb F_{101}[t]$、$B=S[y]/(y^2-F)$ とする。中心は前々節の四切断を
四位置の共通分母で多項式化する。各表示と始点0,1について、
\path{computations/verify_rank_fifteen_all_ramification.py} は
定理\ref{thm:rank-fourteen-geometric-reduction}の双対経路の計算を行う。
$B$ を $S$ 上の階数2の加群として、各段で適用する関係の行列 $A$、
$\ker A$ の基底と、この基底での $y$ の作用を求める。
独立な像の列からなる行列の行加群を取り、その正方基底行列の行列式を使う。
この行列式が単元となる局所環では像が直和因子になり、核が剰余体への
基底変換と交換する。従って前段からの矢の写像も含めて同じ計算が特殊化する。

四始点のうち二つと二表示について、長さ2から13までの全48文字ブロックをそれぞれ調べた。
一般次元は標準中心商の各次数の次元の二倍である。$F(t)\ne0$ の剰余点では
二つの $y$ の固有部分を考える。一般には $y^2-F$ が既約で両部分の次元が等しく、
核が基底変換と交換する領域ではエタールな二次被覆上で各部分が同じ階数を保つ。
例外因子を次の表に示す。中央二列は各始点の行列式の相異なる既約因子数、
最終列は両始点の和集合から境界と一次因子を除いた後の個数である。
\[
\begin{array}{c|rr|r}
 \text{表示}&\text{始点0}&\text{始点1}&\text{残る因子}\\ \hline
 1&29&29&8\\
 2&26&26&7
\end{array}
\]
各残余因子の根で四次中心を取り直す。$y$ が剰余体で分解するときは両根、
非分解なら二次拡大上の一根で、全四始点の商次元が標準列に一致する。
後者は相方と共役である。二位置移動の像で例外になる点も、この検査結果を
移すことで扱える。従ってどの通常な還元でも上界が持ち上がる。
前節の境界円板と合わせ、代数的な $t$ は数体上の付値を取れば全て尽くされる。
超越的な $t$ は一般点の体拡大である。中心の持上げと自由性は前節と同じである。
最後に定理\ref{thm:rank-fifteen-branch-center}で中心次元が全点で分かっている。
系\ref{cor:rank-fifteen-near-collision}の有限射の議論を分岐の全窓にも適用する。
\end{proof}

\begin{proposition}\label{prop:rank-fifteen-pencil}
二つの階数低下被覆の各々において、$\xi=zm^2/n^4$ は通常域で可逆な関数である。
$\xi=-1$ が階数14、$\xi=1$ が被約分岐曲線である。
他の各 $\xi$ のファイバーは幾何学的に既約で、滑らかな完備化の種数は
$\xi^2+\xi+1=0$ で3、その他で5となる。
二位置移動は各完備曲線上で不動点を持たず、商の種数はそれぞれ2と3である。
二つの種数3の値では両表示の一般窓が標準Koszul左右AS正則である。
\end{proposition}
\begin{proof}
幾何学的な値 $\xi=c^2\ne-1$ を固定し、$r=cn^2/m$、
$w=\lambda r$、$K=rk$ と置く。すると $r^2=z$ で、
\[
 n^2=(K^2-1)/(1+c^2),\quad
 (r+r^{-1})/2=K+cn^2,\quad (w+w^{-1})/2=K-cn^2.
\]
逆にこれらの式から $z,\lambda,n$ を復元でき、元の双二次式を満たす。
$P(n)=0$ から $zk^2=1$ は $\xi=-1$ と同値で、判別式は
$n^4(1-\xi)^2$ である。$s=c+c^{-1}$ とすると $K$ の射影直線上の関数体は
次の三つの平方根で生成される。
\[
 \sqrt{K^2-1},\qquad \sqrt{(K+s)^2-1},\qquad \sqrt{(K-s)^2-1}.
\]
実際 $(K\pm cn^2)^2-1= n^2 c^2 ((K\pm s)^2-1)/(1+c^2)$ である。
$r,w$ の二つのトレースは $K$ の非定数式で、恒等的に等しくも符号違いにもならない。
このため通常条件は各曲線の有限集合だけを除く。

三つの二次式の根対は $\{-1,1\}$ とその $\pm s$ 平行移動である。
$s\ne0$ で互いに重なるのは $s=\pm1,\pm2$ のときに限る。
これらの場合も根対を頂点間の辺と見ると閉路がなく、どの非空な積にも奇数位数の零点が残る。
従って拡大の次数は8である。分岐は各幾何根で指数2、無限遠で不分岐となる。
相異なる根の数は一般に6、$s=\pm1$ で5、$s=\pm2$ で4であり、
Riemann--Hurwitzから種数は5,3,1である。$s=\pm1$ が
$\xi^2+\xi+1=0$、$s=\pm2$ が $\xi=1$ に当たる。
二位置移動は $n$ のみを反転し $r,w$ を保つので、三平方根を同時に反転する。
三つの根対に共通点がないため、この元はどの惰性群にも属さない。

$s=1$、$K=3$ で $c^2-c+1=0$ の両根を取る。
139において二根と各 $n,r,w$ の二択は全て単純に分解する。
\path{computations/verify_rank_fifteen_pencil.py} は両表示の合計32窓で
通常条件、反対階数15、中心の文字次元と全四始点の長さ13までの商を確認する。
単純根の持上げと中心の開条件から、各種数3曲線の一般点も所要の正則性を持つ。
\end{proof}

\begin{corollary}\label{cor:rank-fifteen-octic}
両被覆の二つの種数3の特殊値は、幾何学的にはすべて同じ超楕円曲線
\[
 Y^2=t^8+14t^4+1
\]
を持つ。二位置移動は $(t,Y)\mapsto(-1/t,-Y/t^4)$ で、その商の種数2曲線は
$Z^2=(X^2+4)(X^4+4X^2+16)$ である。
\end{corollary}
\begin{proof}
$c^2-c+1=0$、$j^2=c$ を取り $H=jn$ とする。
$H^2=K^2-1$、$A^2=K(K+2)$、$B^2=K(K-2)$ と置くと、
$r=K+H^2+HA$、$w=K-H^2+HB$ が前命題の逆表示になる。
$U=A/K,V=B/K$ は $U^2+V^2=2$ を満たす。$t=(V-1)/(U-1)$ により
\[
 U=\frac{t^2-2t-1}{t^2+1},\quad V=\frac{1-2t-t^2}{t^2+1},\quad
 D=2t(t^2-1),\quad K=-\frac{(t^2+1)^2}{D}.
\]
$Y=DH$ で $D^2(K^2-1)=t^8+14t^4+1$ となり、各式は関数体で逆に解ける。
二位置移動は $H,A,B$ の全符号を反転するので $t'=(V+1)/(U+1)=-1/t$ である。
$D(t')=D(t)/t^4$ から $Y'$ が従う。$X=t-1/t$ と $Z=(t+1/t)Y/t^2$
は不変であり、上記の式を満たす。$t$ を添加すれば $Y$ が戻るので、これらが
不変体を生成する。両モデルの右辺は平方自由である。
\end{proof}

\begin{corollary}\label{cor:rank-fifteen-pencil-boundary}
$c^2\ne-1$ を固定した完備曲線から通常窓を得るために取り除く幾何点の個数は、
$s^2=1$ で44、$s^2\in\{4,-4,9/4\}$ で48、それ以外で56である。
\end{corollary}
\begin{proof}
$R=K+(K^2-1)/s$、$W=K-(K^2-1)/s$ は $r,w$ の半トレースである。
$K$ が有限なら $r,w$ は単元である。$r^4=1$ は $R=0,\pm1$、
$w^2=1$ は $W=\pm1$ と同値となる。残る $w=\pm r,\pm r^{-1}$ は
$W=\pm R$ と同値で $K=0,\pm1$ を与える。
従って境界は、前命題の全分岐値、$K=\infty,0$、$K^2+sK-1=0$
の上の全点である。無限遠では半トレースに極がある。
最後の二次式は $s^2=-4$ で重根を持つ。$\pm1,-s\pm1$ では零にならず、
$s\pm1$ における二つの値の積は $s^2(4s^2-9)$ である。$K=0$ が既に分岐値に
含まれるのは $s^2=1$ の場合である。各分岐値で4点、他の各値で8点を数えればよい。
\end{proof}
